\section{Conclusion}
\label{sec:conclusion}

We presented a cross-resolution supervision method for improving the released
M7 papyrus-surface segmenter without changing its deployed architecture.
A fine teacher exposes local inconsistencies as soft occupancy; a masked
function anchor holds the student to the incumbent wherever the teacher is
silent or confidently agrees, and releases it everywhere else.
Under matched inference the resulting student improves on released M7 by
$\fzGain$ macro-scroll Dice on a frozen benchmark, and by more with an
inference-time augmentation that leaves the weights untouched.
On the same rows it also leads M7 as the organisers publish it and HercUNet v0,
and its solid-mass share does not rise with compression as the published M7's
does.
We release the shipped weights as \releaseName.

The larger lesson of this work is subtractive.
We began with seven weighted training terms, each a defensible answer to an
observed failure, and finished with three.
What removed the others was not a better idea but a better measurement: a
corpus wide enough for a shape prior to be worth less than data, a schedule run
to its declared length, and the removal of a parameter-space constraint that
had been quietly making every ablation uninformative.
The terms we discarded were not obviously wrong when we added them, and they
were not shown to be wrong by argument.
They were shown to be wrong by measuring them in a regime where they could have
been right.

The same design principle survives downstream: the 2D fitter may propose a
connection, but radial, merger, cross-sheet, crossing, and persistence gates
force it to answer a sequence of geometric objections before adding pixels.
A useful surface model must learn to grow longitudinal structure without buying
that recall through thickness or cross-wrap mergers, and the honest way to know
whether a given mechanism does that is to give it a fair chance to fail.
